\documentclass[16pt]{article}
\usepackage{amsmath}


\begin{document}
\section*{Exercice 01:}
\textbf{Formule 1}  On a $12 \div 3 = 4$ signifie aussi $3 \times 4 = 12$.\\
\textbf{Formule 2}  L'égalité $x = 2y$ est équivalente à $y = x/2$.\\
\textbf{Formule 3}  L'égalité $x = 2y$ est équivalente à $y = \frac{1}{2}x$.\\
\textbf{Formule 4}  La fonction qui à $x$ associe $x^2$.\\
\textbf{Formule 5}  On considère la fonction $f$ définie sur \textbf{R}. $$f : x\mapsto e^{2x+1}$$
\textbf{Formule 6}  $$\frac{\frac{x^2}{y^2}}{\frac{y^2}{x^2}} = \frac{x^4}{y^4}$$
\textbf{Formule 7}  $$\dfrac{\dfrac{x^2}{y^2}}{\dfrac{y^2}{x^2}} = \frac{x^4}{y^4}$$
\textbf{Formule 8}
Dans le repère $(O;\vec{i};\vec{j})$,les sous-espaces $\mathcal{F}$ et $\mathcal{G}$
vérifient $\forall (\vec{u},\vec{v}) \in \mathcal{F} \times \mathcal{G}, \ \vec{u}\cdot\vec{v} = 0.$\\
\textbf{Formule 9}$$\frac{\sqrt{x+1}}{y+1} \neq \sqrt{\frac{x+1}{y+1}}$$
\textbf{Formule 10} $$1 + \frac{1}{1 + \dfrac{1}{1 + \dfrac{1}{1 + \dfrac{1}{1}}}} = 1,6$$

\section*{Exercice 02:}
\textbf{Formule 1}  On a $\sin \frac{\pi}{2} = 1$ et $\sin \frac{\pi}{4} = \frac{\sqrt{2}}{2}$ \\
\textbf{Formule 2}
\[ \sum_{i,j\in I\times J} i+j = \sum_{i\in I}\left(\sum_{j\in J} i+j\right) \]
 \\
\textbf{Formule 3}
\[ \sum_{p\in\mathcal{P}}\frac{1}{p} \qquad\text{ou}\qquad
   \sum_{p\text{ premier}}\frac{1}{p} \]
 \\
\textbf{Formule 4} \quad
Par définition, \(nx=\overbrace{x+\cdots+x}^{n\text{ fois}}\).
 \\
\textbf{Formule 5} \quad
On a : \(nx \overset{\text{déf}}{=} \overbrace{x+\cdots+x}^{n\text{ fois}}\).
 
\medskip
\textbf{Formule 6}\[ f : \begin{array}{c} \mathbf{R}\to\mathbf{R} \\ x\mapsto e^x \end{array} \]
On verra lors de la correction comment obtenir le résultat(plus classique) :\[ f : \begin{array}{rcl} \mathbf{R} & \to & \mathbf{R} \\
                          x & \mapsto & e^x \end{array} \]\\
\textbf{Formule 7}
\[ \iint_{\mathcal{D}} z^2\,\mathrm{d}z \]
La difficulté, ici, est de gèrer correctement les espacements.\\
\textbf{Formule 8}
\[ \underbrace{\overbrace{a,\ldots,a}^{k},\overbrace{c,\ldots,c}^{l}}_{k+l\text{ éléments}} \]\\
\textbf{Formule 9}
\[
\begin{pmatrix}
  \begin{pmatrix} 1 & 0 \\ 0 & 1 \end{pmatrix} &
  \begin{pmatrix} a & b \\ c & d \end{pmatrix} \\
  \begin{pmatrix} d & -b \\ -c & a \end{pmatrix} &
  \begin{pmatrix} 0 & -1 \\ 1 & 0 \end{pmatrix}
\end{pmatrix}
\]

\end{document}